%% 60-impact.tex — Impact and Adoption

\section{Impact and Adoption}
\label{sec:impact}

\subsection{Access to the DDB Corpus}

GEMEA is the first open, SPARQL-accessible knowledge graph over the DDB corpus, making the digitized collection queryable by SPARQL, dereferenceable as Linked Data, and accessible to AI agents via MCP.

\subsection{FAIR Compliance}

GEMEA satisfies the FAIR data principles~\cite{wilkinson2016fair}:

\begin{itemize}
    \item \textbf{Findable.} The dataset is publicly available for download at \url{https://gemea.ise.fiz-karlsruhe.de/downloads}.
    \item \textbf{Accessible.} Data is retrievable via a public SPARQL endpoint and through the SHMARQL Linked Data browser; entity URIs support content negotiation (HTML, Turtle, JSON-LD).
    \item \textbf{Interoperable.} Standard RDF serializations (N-Triples, Turtle, JSON-LD), established vocabularies (EDM, RDA, FRBR, Dublin Core, GeoNames), and GND authority identifiers for agents and works are used throughout.
    \item \textbf{Reusable.} Source code, pipeline scripts, and MCP setup are released at \url{https://github.com/anntanp/gemea/} under the MIT licence. Individual object rights are as declared by the contributing institution (typically CC~1.0) and are queryable via \texttt{dcterms:rights} in the graph.
\end{itemize}

\subsection{Target Communities}

\paragraph{Digital humanities and library science.}
German cultural heritage research has lacked a machine-readable, graph-accessible corpus at this scale.
GEMEA enables quantitative analysis previously infeasible without custom DDB API scraping: cross-sector co-occurrence patterns, temporal distribution of holdings, and geographic coverage.

\paragraph{Semantic Web and KG research.}
A public SPARQL endpoint, named graph provenance, \textbf{MOCHO} ontology alignment, and MCP agentic access make GEMEA a testbed for KG construction methods, ontology evaluation, entity linking, and LLM + KG research (Section~\ref{sec:usage}).

\paragraph{Knowledge graph embedding research.}
GEMEA's \textbf{MOCHO}-aligned representation of bibliographic titles opens a concrete empirical question that remains partially explored in the KGE-with-literals literature~\cite{kristiadi2019incorporating,xie2016representation,wang2021kepler,blum2022literaltransformation}: whether entity descriptions yield better embeddings when modeled as a single rich literal or as multiple typed relations with shorter, semantically structured values.
Existing approaches incorporate literals either by embedding them directly (e.g., LiteralE~\cite{kristiadi2019incorporating}, DKRL~\cite{xie2016representation}, KEPLER~\cite{wang2021kepler}) or by transforming them into graph structure~\cite{blum2022literaltransformation}.  However, these methods treat the representation of literals as a modeling choice fixed at preprocessing time, and do not systematically compare alternative semantic decompositions of the same information. In particular, prior work focuses on syntactic or structural transformations of literals (e.g., converting values into entities), rather than semantically grounded decompositions aligned with domain ontologies.
GEMEA makes this comparison tractable: the \textbf{MOCHO}-aligned subgraph exposes both the decomposed RDA title structure and the source \texttt{dc:title} literals within the same graph, enabling controlled ablation studies at scale on a real-world cultural heritage corpus.


\paragraph{Cultural heritage institutions.}
GEMEA provides a self-hostable reference deployment that any cultural heritage aggregator can adapt: the JSON ingest + \texttt{mocho} + QLever stack, with SHMARQL as a Linked Data browser, MCP/MCPO integration to commercial AI assistants and other self-hosted setups (Open WebUI + Ollama). The full stack is documented, containerized, and deployable via Docker Compose (Section~\ref{sec:sustainability}).
The more than 500 institutions that have contributed data to the DDB (out of over 5,000 registered) can query their own holdings in graph form through GEMEA, enabling metadata introspection, coverage analysis, and targeted enrichment that feeds back into the aggregator.
